\section{Experimental Details}
\label{xsec:experimental_details}


% 图片merge的规则：
\subsection{Merging Rules of Images Input}
\label{xsubsec:rules_of_images_merging}
As discussed in \Cref{subsec: ablation_of_input_paradigms}, we design a \textbf{merge} paradigm for the evaluation of current LLMs/LVLMs.

\noindent\texttt{\#Columns\_merged = 2 if total\_pages/30 <= 4 else 3}

\begin{table}[htbp]
    \centering
    \resizebox{0.5\textwidth}{!}{
    \begin{tabular}{l|c|c}
        \toprule
        \textbf{Total\_pages} & \textbf{\#Columns\_merged} & \textbf{\#Images\_merged} \\
        \midrule
        50<x<=60 & 2 & x/2 \\
        60<x<=90 & 2 & x/3 \\
        90<x<=120 & 2 & x/4 \\
        120<x<=150 & 3 & x/5 \\
        \bottomrule
    \end{tabular}
    }
    \caption{Merging rules in the ablation experiment of input paradigms. \textbf{\#Columns\_merged}: the number of columns in the sub-image array in the merged image; \textbf{\#Images\_merged}: the number of new images after merged.}
    \label{tab:merging_rules}
    \vspace{-5mm}
\end{table}


% 输入图片30页的选择规则：
\subsection{Selection Rules of Images Input}
\label{xsubsec:rules_of_images_selection}
As discussed in \Cref{subsec: experimental_setup}, we design a \textbf{cut-off} paradigm for the evaluation of current LLMs/LVLMs. We provide pseudo-code below to express the selection rules.

\begin{tcolorbox}[colback=gray!5,left=0.8mm,right=0.8mm]
\small
\textcolor{gray}{\# All page ids mentioned are based on the order in the raw document.} \\
\textcolor{gray}{\# context\_start\_end: page start/end id of context in qa generation} \\
\textcolor{gray}{\# num\_of\_images\_input: the number of input images, 30 images in cut-off paradigm} \\
\textcolor{gray}{\# total\_pages: total pages in the raw document} \\
\textcolor{gray}{\# img\_start, img\_end: input page start/end id} \\
\raggedright
\texttt{raw\_start\_page, raw\_end\_page = context\_start\_end} \\

\texttt{raw\_pages\_len = raw\_end\_page - raw\_start\_page} \\

\texttt{img\_start = max(0, raw\_start\_page -}  \\
\texttt{\quad \quad (num\_of\_images\_input - raw\_pages\_len)//2} \\
% \texttt{\quad // 2)} \\

\texttt{img\_end = img\_start + num\_of\_images\_input} \\

\texttt{if img\_end >= total\_pages:} \\

\texttt{\quad \quad img\_end = total\_pages} \\

\texttt{\quad \quad img\_start = max(0, } \\
%img\_end-num\_of\_images\_input)}\\
\texttt{\quad \quad img\_end - num\_of\_images\_input)} \\

\texttt{return img\_start, img\_end}
        
\label{text:selection_rules}
\end{tcolorbox}


% 1. 自由生成回复
\subsection{Prompt for Response Generation}
\label{xsubsec:prompt_of_response_generation}

% \lstinputlisting[style=prompt]{./prompts/response_generation.md}

\begin{tcolorbox}[colback=gray!5]
\small
\textcolor{blue}{<document\_images>} \\

You are an expert in visual document question-answering, please answer our questions based
on the given doc images.\\

Following is our question:\\
<question\_start>\textcolor{blue}{<question>}</question\_end>

\label{text:prompt_for_response_generation}
\end{tcolorbox}


% 2. 提炼答案
\subsection{Prompt for Answer Extraction}
\label{xsubsec:prompt_of_answer_extract}

Prompt for answer extraction is displayed in \Cref{fig:answer_extraction}. Based on the template given in MMLongBench-Doc~\cite{MMLongBench-Doc}, we make some modifications, which are marked in \textcolor{blue}{blue}.


\begin{figure*}[htbp]
\begin{tcolorbox}[title=The answer extracting prompt used to make long response concise during evaluation,colback=gray!5,colframe=gray!75!gray]
% \begin{tcolorbox}[breakable, title=Porceeding for task: Place the red cube 3cm above the brown drawer,colback=gray!5,colframe=nmgray!75!black,before skip=1pt, after skip=1pt,fontupper=\linespread{0.8}\selectfont]
{\footnotesize

% \vspace{0.2cm}

\vspace{0.1cm}
\textbf{[Task Description]}

Given the question and analysis, you are tasked to extract answers with required formats from the free-form analysis. 

\vspace{0.1cm}

Your extracted answers should be one of the following formats: (1) Integer, (2) Float, (3) String and (4) List. If you find the analysis the question can not be answered from the given documents, type "Not answerable". Exception: If the analysis only tells you that it can not read/understand the images or documents, type "Fail to answer". 

\vspace{0.1cm}

\textcolor{blue}{Please use <concise\_answer> and <\/concise\_answer> tokens at the start and end of the extracted answer. For example, if the extracted answer is number 3, the format is <concise\_answer>3<\/concise\_answer>.}

\vspace{0.1cm}

\textcolor{blue}{Please use <answer\_format> and <\/answer\_format> tokens at the start and end of the answer format. For example, if the answer format is List, the format is <answer\_format>List<\/answer\_format>.}


\vspace{0.1cm}

Please make your response as concise as possible. Also note that your response should be formatted as below: 


\vspace{0.2cm}


\textcolor{blue}{Extracted answer: <concise\_answer>[answer]<\/concise\_answer>}

\textcolor{blue}{Answer format: <answer\_format>[answer format]<\/answer\_format>}

\vspace{0.2cm}

\textbf{[Few Shots]}

Please read the following example, then extract the answer from the model response and type it at the end of the prompt. 

\vspace{0.2cm}

--- Example 1:

\textbf{Question}: List the primary questions asked about the services in this report.

\textbf{Analysis}: The primary questions asked about the services in the report for The Limes Residential Home are: 
% \begin{enumerate}[itemsep=0.01ex]
%     \item Is the service safe?
%     \item Is the service effective?
%     \item  Is the service caring?
%     \item Is the service responsive?
%     \item Is the service well-led? 
% \end{enumerate}

1. Is the service safe?

2. Is the service effective?

3. Is the service caring?

4. Is the service responsive?

5. Is the service well-led? 

\textbf{Extracted answer}: <concise\_answer>['Is the servife safe?', 'Is the service effective', 'Is the serve caring?', 'Is the service responsive?', 'Is the service well-led?']<\/concise\_answer>

\textbf{Answer format}: <answer\_format>List<\/answer\_format>

\vspace{0.2cm}

--- Example 2:

\textbf{Question}: How many regulations of the HSCA 2008 are breached in all according to this report?

\textbf{Analysis}: According to the report, the provider breached 10 Health and Social Care Act 2008 (Regulated Activities) Regulations in total. Here are the specifics:

1. Regulation 13: Safeguarding service users from abuse and improper treatment

2. Regulation 12: Safe care and treatment

3. Regulation 18: Staffing

4. Regulation 11: Need for consent

5. Regulation 10: Dignity and respect

6. Regulation 9: Person-centred care

7. Regulation 17: Good governance

8. Regulation 18 (CQC Registration Regulations 2009): Notification of other incidents

9. Regulation 18: Failure to maintain an accurate and up-to-date care plan

10. Regulation 11: Failure to implement the Mental Capacity Act 2005 code of practice effectively

These breaches involve issues concerning staffing, safeguarding, medicines management, dignity and respect, consent, care planning, governance, and failure to notify the CQC of incidents.

\textbf{Extracted answer}: <concise\_answer>10<\/concise\_answer>

\textbf{Answer format}: <answer\_format>Integer<\/answer\_format>

\vspace{0.2cm}

--- Example 3:

\textbf{Question}: According to the survey that is the percentage of Chinese who are paying more or about the same attention to politics after Trump’s election?

\textbf{Analysis}: The survey provided does not specify the percentage of Chinese individuals specifically who are paying more or about the same attention to politics after Trump’s election. The report focuses primarily on American demographics and does not include specific details about the Chinese population in relation to this question. If you need information about a different demographic or a summary of the findings from the American demographic, I can certainly help with that!

\textbf{Extracted answer}: <concise\_answer>Not answerable<\/concise\_answer>

\textbf{Answer format}: <answer\_format>None<\/answer\_format>

\vspace{0.2cm}

--- Example 4:

\textbf{Question}: How many quotations from male respondent over 50 years old are included in this report?

\textbf{Analysis}: The image you’ve provided appears to be a screenshot of a document with multiple charts. However, the text is too small and blurry to read accurately. If you can provide a clearer image or more context, I might be able to help you with your question.

\textbf{Extracted answer}: <concise\_answer>Fail to answer<\/concise\_answer>

\textbf{Answer format}: <answer\_format>None<\/answer\_format>


}
\end{tcolorbox}
\vspace{-0.2cm}
\caption{Prompt for answer extraction during evaluation.}
\label{fig:answer_extraction}
\end{figure*}


% 3. 计算分数
\subsection{Scoring Rules}
\label{xsubsec:scoring_rules_details}

Following MATHVISTA~\cite{MATHVISTA}, we evaluate the model's responses by scoring the extracted answers against the reference answers. Following MMLongBench-Doc~\cite{MMLongBench-Doc}, our scorer is rule-based and employs different strategies according to the format of the reference answer. We detail its rules as below:

\noindent\textbf{String:} We firstly use a series of regular expressions to determine whether the answers require exact matching (\eg, telephone numbers, email addresses, website addresses, file names, times, dates, \etc.) If an exact match is needed, we perform a straightforward string comparison and score the answer either 0 or 1. Otherwise, we calculate the ANLS (Average Normalized Levenshtein Similarity) with a pre-defined threshold ($\tau = 0.5$).

\noindent\textbf{Integer:} We perform an exact match comparison and score the answer either 0 or 1.

\noindent\textbf{Float:} We view the prediction and reference answers as equal if they fall within a 1\% relative tolerance. 

\noindent\textbf{List:} Compared with MMLongBench-Doc, we adopt a relatively soft rule for scoring answers in list format: (1) If the prediction does not have the same number of elements as the reference, it incurs a length-dependent penalty instead of receiving a score of 0, which we think more reasonable. (2) The score of models on single element of the reference list is the highest one among the scores which are calculated and combined between the element and each one of the prediction list. Compared with MMLongBench-Doc, we assume that the sorting positions of the two lists are not always one-to-one corresponding, allowing more errors, and our rules are gentler and more tolerant. We use pseudo-code below to express the scoring rules in MMLongBench-Doc and our \benchName, respectively. The element-wise scoring strategies is determined by the formats of elements (\ie, string, integer or float).

\begin{tcolorbox}[colback=gray!5]
\small
\textcolor{gray}{\# MMLongBench-Doc} \\
pred\_list, ref\_list = sorted(pred\_list), sorted(ref\_list) \\
Score(pred\_list, ref\_list) = min([score(pred, ref) for pred, ref in zip(pred\_list, ref\_list)]) \\

\vspace{0.5em}
\textcolor{gray}{\# LongDocURL} \\
pred\_list, ref\_list = sorted(pred\_list), sorted(ref\_list) \\
greedy\_scores\_list = [ \\
\setlength{\parindent}{2em} 
\hspace{1em} max([score(pred, ref) for pred in pred\_list]) for ref in ref\_list \\
] \\
Score(ref\_list, pred\_list) = \\ 
\hspace{1em} sum(greedy\_scores\_list) / len(ref\_list) * min(1, len(ref\_list) / len(pred\_list)) ** 0.5 \\

\label{text:scoring_rules_for_list}
\end{tcolorbox}

The part of the rule description that differs from MMLongBench-Doc is mainly in the \textbf{List} section.


\subsection{Fine-Grained Evaluation Results}
\label{xsubsec:appendix_multi_dimensional_results}
Detailed results are presented in \Cref{tab:main_results_2}. 

\noindent Related analysis is in \Cref{subsubsec:multi-dimensional_analysis}.

% \vfill
\begin{table*}[htbp]
    \centering
    \resizebox{\textwidth}{!}{
    \begin{tabular}{lc|ccc|cccc|ccc|c}
    % \begin{tabular}{lcc|ccc|cccc|ccc|c}
        \toprule
        \multirow{3}{*}{\textbf{Model}} & \multirow{3}{*}{\textbf{Size}} & \multicolumn{3}{c|}{\textbf{Task Type}} & \multicolumn{4}{c|}{\textbf{Evidence Element}} & \multicolumn{3}{c|}{\textbf{Page/Element}}& \multirow{2}{*}{\textbf{Total}} \\

        & & U & R & L & TXT & LAY & FIG & TAB & SP & MP & CE & \\
        & & 1243 & 387 & 695 & 994 & 779 & 556 & 871 & 1093 & 1230 & 862 & 2325 \\
        
        \midrule
        
        \multicolumn{13}{c}{\textit{OCR (PyMuPDF\footnotemark[4]) + Large Language Models (LLMs)} } \\
        
        \midrule
        
        \textcolor{gray}{\textit{Open-source Models}} \\
        LLaVA-Next-Interleave & 7B & 26.4 & 17.2 & 5.8 & 26.0 & 13.8 & 15.4 & 14.0 & 14.4 & 22.6 & 16.0 & 18.7 \\
        LLaVA-Next-Interleave-DPO & 7B & 28.7 & 16.4 & 6.4 & 27.4 & 15.7 & 17.7 & 14.2 & 15.1 & 24.3 & 17.5 & 20.0 \\
        LLaVA-OneVision & 7B & 31.0 & 20.2 & 11.2 & 32.6 & 19.2 & 22.1 & 15.1 & 17.5 & 28.5 & 21.0 & 23.3 \\
        LLaVA-OneVision-Chat & 7B & 31.7 & 20.9 & 14.0 & 33.1 & 19.8 & 25.3 & 16.6 & 20.0 & 28.8 & 22.2 & 24.6 \\
        Qwen2-VL & 7B & 31.4 & 22.5 & 14.9 & 33.2 & 20.6 & 25.0 & 17.2 & 20.7 & 28.7 & 22.6 & 25.0 \\
        Qwen2.5-Instruct & 7B & 29.6 & 23.7 & 20.5 & 31.4 & 27.0 & 24.3 & 18.2 & 21.5 & 29.7 & 24.2 & 25.9 \\
        Qwen2.5-Instruct & 14B & 31.3 & 23.9 & 24.2 & 33.3 & 29.3 & 24.8 & 21.6 & 23.5 & 31.8 & 26.5 & 27.9 \\
        Qwen2.5-Instruct & 32B & 30.0 & 25.7 & 29.5 & 31.9 & 30.4 & 28.9 & 25.8 & 26.4 & 31.6 & 29.6 & 29.2 \\
        Qwen2.5-Instruct & 72B & 33.9 & 27.2 & 34.2 & 37.4 & 33.8 & 34.0 & 28.6 & 29.3 & 36.0 & 35.3 & 32.9 \\

        \textcolor{gray}{\textit{Proprietary Models}} \\
        Qwen-Max & - & 32.1 & 24.5 & 34.0 & 34.6 & 32.8 & 32.6 & 27.7 & 28.3 & 34.1 & 33.2 & 31.4 \\
        Gemini-1.5-Pro & - & 33.3 & 26.7 & 32.8 & 35.5 & 34.7 & 32.4 & 26.9 & 28.3 & 35.4 & 33.1 & 32.0 \\
        Qwen-VL-Max & - & \colorbox{firstBest}{\textbf{37.7}} & 30.0 & 27.3 & \colorbox{firstBest}{\textbf{39.6}} & 31.2 & 35.1 & 27.9 & 31.1 & 35.2 & 34.2 & 33.3 \\
        GPT-4o & - & 35.4 & \colorbox{firstBest}{\textbf{28.3}} & 37.2 & 37.6 & 36.3 & 36.3 & 30.2 & 31.9 & 37.2 & 35.8 & 34.7 \\
        O1-preview & - & 35.6 & 31.2 & \colorbox{firstBest}{\textbf{38.6}} & 37.4 & \colorbox{firstBest}{\textbf{37.1}} & \colorbox{firstBest}{\textbf{37.3}} & \colorbox{firstBest}{\textbf{34.6}} & \colorbox{firstBest}{\textbf{33.5}} & \colorbox{firstBest}{\textbf{37.8}} & \colorbox{firstBest}{\textbf{38.6}} & \colorbox{firstBest}{\textbf{35.8}} \\

        \midrule
        
        \multicolumn{13}{c}{\textit{Large Vision Language Models (LVLMs)} } \\
        
        \midrule
        
        \textcolor{gray}{\textit{Open-source Models}} \\
        InternLM-XC2.5 & 7B & 3.6 & 1.8 & 0.7 & 3.9 & 2.2 & 1.9 & 1.2 & 1.9 & 2.9 & 2.3 & 2.4 \\
        mPLUG-DocOwl2 & 7B & 7.7 & 3.8 & 1.8 & 8.1 & 4.7 & 5.5 & 3.1 & 3.7 & 6.7 & 6.1 & 5.3 \\
        Pixtral & 12B & 7.9 & 3.7 & 2.4 & 8.8 & 4.8 & 4.3 & 3.1 & 4.8 & 6.2 & 5.3 & 5.6 \\
        Llama-3.2 & 11B & 12.9 & 9.4 & 2.7 & 11.8 & 6.9 & 8.7 & 6.3 & 7.9 & 10.3 & 6.8 & 9.2 \\
        LLaVA-Next-Interleave & 7B & 20.2 & 13.0 & 3.8 & 21.8 & 10.8 & 12.1 & 6.7 & 10.0 & 17.8 & 11.3 & 14.1 \\
        LLaVA-Next-Interleave-DPO & 7B & 21.6 & 13.9 & 7.6 & 22.5 & 13.9 & 15.4 & 8.7 & 12.1 & 19.8 & 13.5 & 16.2 \\
        LLaVA-OneVision & 7B & 28.1 & 16.5 & 14.1 & 30.8 & 23.9 & 17.9 & 11.6 & 16.5 & 26.8 & 20.8 & 22.0 \\
        LLaVA-OneVision-Chat & 7B & 30.5 & 19.0 & 18.7 & 32.2 & 26.5 & 24.4 & 15.4 & 19.8 & 29.7 & 24.2 & 25.0 \\
        Qwen2-VL & 7B & 36.9 & 24.8 & 22.6 & 37.7 & 29.7 & 28.6 & 23.7 & 27.2 & 33.6 & 29.9 & 30.6 \\
        
        \textcolor{gray}{\textit{Proprietary Models}} \\
        Qwen-VL-Max & - & 58.8 & 43.9 & 36.0 & 58.0 & 40.2 & 52.3 & 44.6 & 51.6 & 47.6 & 48.0 & 49.5 \\
        Gemini-1.5-Pro & - & 55.7 & 43.4 & 46.4 & 58.7 & 50.4 & 50.0 & 41.8 & 48.7 & 52.8 & 49.9 & 50.9 \\
        GPT-4o & - & \colorbox{firstBest}{\textbf{68.6}} & \colorbox{firstBest}{\textbf{59.9}} & \colorbox{firstBest}{\textbf{59.6}} & \colorbox{firstBest}{\textbf{70.7}} & \colorbox{firstBest}{\textbf{60.0}} & \colorbox{firstBest}{\textbf{67.4}} & \colorbox{firstBest}{\textbf{60.3}} & \colorbox{firstBest}{\textbf{65.8}} & \colorbox{firstBest}{\textbf{63.2}} & \colorbox{firstBest}{\textbf{65.4}} & \colorbox{firstBest}{\textbf{64.5}} \\

        \midrule
        \multicolumn{13}{c}{\textit{Human Baseline} } \\
        \midrule
        Human Experts & - & - & - & - & - & - & - & - & - & - & - & 84.8 \\
        
        \bottomrule
    \end{tabular}
    }
    \caption{Generalized accuracy scores(0\textasciitilde1) on \benchName. There are 3 types of tasks: (U)nderstanding, (R)easoning, and (L)ocalization. There are 4 types of evidence elements: pure text(TXT), layout(LAY), chart \& image(FIG), and table(TAB). There are 3 types of evidence pages/elements: single-page(SP), multi-page(MP), and cross-element(CE). Cross-element means at least two element types in the evidence(\eg, chart and table). The highest scores among models in each section are highlighted in \colorbox{firstBest}{\textbf{green}}.}
    \label{tab:main_results_2}
    \vspace{-5mm}
\end{table*}
